\begin{table}[pos=htbp]
  \centering
  \caption{Observed joint mesh/tolerance refinement response for the coupled nine-species Case 3C calculation. Both calculations use the same equations and parameters. The iteration rows show the printed maximum normalized RMS residual after each mesh iteration; final residuals below retain additional precision.}
  \label{tab:reactive-numerical-verification}
  \small
  \renewcommand{\arraystretch}{1.12}
  \begin{tabularx}{\textwidth}{@{}Xrr@{}}
    \toprule
    Quantity & Coarser calculation & Refined calculation \\
    \midrule
    Initial / final mesh nodes & 21 / 22 & 41 / 42 \\
    Collocation tolerance & 0.5 & 0.05 \\
    Mesh iterations & 2 & 2 \\
    Maximum RMS residual, iterations 1 $\to$ 2 & $1.64\to0.123$ & $0.0507\to0.0393$ \\
    Final maximum RMS residual & 0.122785 & 0.039322 \\
    Maximum scaled boundary residual & 0 & 0 \\
    Column RHS evaluation batches & 27 & 17 \\
    Jacobian evaluation batches & 6 & 4 \\
    State Jacobian evaluations & 125 & 164 \\
    Liquid equilibrium solves, including profile output & 608 & 551 \\
    \COtwo{} capture (\unit{\percent}) & 94.19959 & 94.20397 \\
    Peak liquid temperature (\unit{K}) & 347.9682 & 347.9399 \\
    Net-enthalpy-flow range, 101 positions (\unit{W}) & 286.440 & 267.572 \\
    \bottomrule
  \end{tabularx}
  \par\smallskip
  \begin{minipage}{\textwidth}\footnotesize
    The continuous equations predict zero net-enthalpy-flow range; the nonzero numerical values are diagnostics, with no asserted energy-error threshold. RHS denotes the right-hand side of the column equations. A batch evaluates all nodes or interval midpoints, not a single axial state. Four Jacobian batches correspond to two node/midpoint pairs. Mesh iterations count outer refinement cycles, not inner Newton steps. Both calculations use boundary tolerance \num{0.001} and a 1000-node limit. Peak memory is reported with the timing study.
  \end{minipage}
\end{table}
